\documentclass[10pt]{article}
\usepackage[T1]{fontenc}
\usepackage[utf8]{inputenc}
\usepackage[french]{babel}
\usepackage{amsmath,amssymb,mathrsfs,mathpazo}
\usepackage{geometry}
\geometry{a4paper,margin=2cm}
\usepackage{graphicx,array,tabularx,tabularray,colortbl,multirow,multicol,booktabs,enumitem,pifont,ulem,fancybox,diagbox,cancel}
\usepackage{pstricks,pstricks-add,pst-plot,pst-tree,pst-node,pst-text,pst-func,pst-math,pst-bspline,pst-3dplot,pst-eucl,pst-solides3d}
\let\pscalculate\undefined
\usepackage{pst-bezier}
\usepackage{tikz,pgfplots,tkz-tab}
\usetikzlibrary{arrows,arrows.meta,patterns,calc,positioning,trees,decorations.markings,shapes.geometric,shapes.misc,shapes.arrows,intersections,decorations.pathmorphing,decorations.pathreplacing,backgrounds,shadows,babel}
\pgfplotsset{compat=1.18}
\usepgfplotslibrary{fillbetween}
\usepackage[np]{numprint}
\usepackage{listings,makecell,pas-tableur,twoopt}
\let\barre\undefined
\usepackage[pstricks]{tablvar}
\definecolor{darkgreen}{RGB}{0,100,0}
\providecommand{\red}{\color{red}}
\providecommand{\blue}{\color{blue}}
\providecommand{\green}{\color{green}}
\providecommand{\black}{\color{black}}
\providecommand{\white}{\color{white}}
\providecommand{\gray}{\color{gray}}
\providecommand{\cyan}{\color{cyan}}
\providecommand{\magenta}{\color{magenta}}
\providecommand{\yellow}{\color{yellow}}
\providecommand{\violet}{\color{violet}}
\providecommand{\orange}{\color{orange}}
\providecommand{\euro}{\texteuro}
\providecommand{\EUR}{\texteuro}
\providecommand{\ieme}{\textsuperscript{e}}
\providecommand{\degres}{\ensuremath{^\circ}}
\providecommand{\degree}{\ensuremath{^\circ}}
\providecommand{\vv}[1]{\overrightarrow{#1}}
\providecommand{\Coord}[1]{#1}
\providecommand{\decosix}{}
\providecommand{\starredbullet}{\ensuremath{\bullet}}
\pagestyle{empty}
\setlength{\parindent}{0pt}
\usepackage{adjustbox}\DeclareUnicodeCharacter{2236}{:}\DeclareUnicodeCharacter{2190}{\ensuremath{\leftarrow}}
\providecommand{\R}{\mathbb{R}}
\providecommand{\N}{\mathbb{N}}
\providecommand{\Z}{\mathbb{Z}}
\providecommand{\Q}{\mathbb{Q}}
\providecommand{\D}{\mathbb{D}}
\providecommand{\C}{\mathbb{C}}
\providecommand{\e}{{\mathrm{e}}}
\providecommand{\vect}[1]{\overrightarrow{#1}}
\providecommand{\vectt}[1]{\overrightarrow{\mathrm{#1}}}
\providecommand{\vvt}[1]{\overrightarrow{\mathrm{#1}}}
\providecommand{\np}[2][]{{#2~#1}}
\providecommand{\vv}[1]{{\overrightarrow{#1}}}
\renewcommand{\d}{{\mathrm{d}}}
\providecommand{\degres}{{^\circ}}
\providecommand{\degre}{{^\circ}}
\providecommand{\blue}{\color{blue}}
\providecommand{\red}{\color{red}}
\providecommand{\enskip}{\quad}
\providecommand{\llbracket}{\lbrack\!\lbrack}
\providecommand{\rrbracket}{\rbrack\!\rbrack}
\providecommand{\euro}{{\text{€}}}
\providecommand{\no}{{\text{n°}}}
\providecommand{\interval}[2]{{[#1;#2]}}
\providecommand{\intervFO}[2]{{[#1;#2[}}
\providecommand{\dotfill}{\ldots}

\providecommand{\ds}{\displaystyle}
\providecommand{\textding}[1]{#1}
\providecommand{\Oij}{\(\mathrm{O};\vec{\imath},\vec{\jmath}\)}
\providecommand{\Oijk}{\(\mathrm{O};\vec{\imath},\vec{\jmath},\vec{k}\)}
\providecommand{\Ouv}{\(\mathrm{O};\vec{u},\vec{v}\)}
\providecommand{\intervFF}[2]{\ensuremath{[#1;#2]}}
\renewcommand{\euro}{\ifmmode\text{\texteuro}\else\texteuro\fi}
\DeclareUnicodeCharacter{00B0}{\ifmmode^{\circ}\else\textdegree\fi}
\DeclareUnicodeCharacter{20AC}{\ifmmode\text{\texteuro}\else\texteuro\fi}
% Énoncés : Éducation nationale. Transcription et corrigés : APMEP, Denis Vergès et contributeurs. Découpage et classement : CoopMaths, Sébastien Lozano, Éric Elter et contributeurs. Adaptation de la présentation : Maths Explorer.
% CC BY-SA : https://coopmaths.fr/www/about/
% https://coopmaths.fr/alea/static/bac/2021/tex/bac_2021_05_ameriquenord_5.tex
\begin{document}
\noindent\textit{APMEP / CoopMaths. CC BY-SA. Présentation adaptée par Maths Explorer.}\par\medskip
\smallskip

\begin{tabular}{|l|}\hline
Principaux domaines abordés : \\ 
\hspace{1cm}$\bullet~~$Fonction logarithme népérien\\
\hspace{1cm}$\bullet~~$Convexité\\ \hline
\end{tabular}

\medskip

Dans le plan muni d'un repère, on considère ci-dessous la courbe $\mathcal{C}_f$ représentative d'une fonction $f$, deux fois dérivable sur l'intervalle $]0~;~ +\infty[$. 

La courbe $\mathcal{C}_f$ admet une tangente horizontale $T$ au point A(1~;~4).

\begin{center}
\psset{unit=1.25cm}
\begin{pspicture*}(-0.5,-0.8)(6.4,4.4)
\psgrid[gridlabels=0pt,subgriddiv=5,gridcolor=gray](0,-1)(9,5)
\psaxes[linewidth=1.25pt]{->}(0,0)(0,-0.8)(6.4,4.4)
\psplot[plotpoints=2000,linewidth=1.25pt,linecolor=red]{0.1}{6.4}{x ln 1 add 4 mul x div}
\psline[linewidth=1.25pt](0,4)(6.4,4)\uput[u](5.8,4){$T$}
\uput[u](5.8,1.9){\red $\mathcal{C}_f$}
\uput[u](1,4){A}
\end{pspicture*}
\end{center}

\medskip

\begin{enumerate}
\item Préciser les valeurs $f(1)$ et $f'(1)$.
\end{enumerate}

On admet que la fonction $f$ est définie pour tout réel $x$ de l'intervalle $]0~;~ +\infty[$ par:

\[f(x) = \dfrac{a + b \ln x}{x} \,\, 
\text{où }\, a \text{ et}\, b \text{ sont deux nombres réels}.\]

\begin{enumerate}[resume]
\item Démontrer que, pour tout réel $x$ strictement positif, on a :

\[f'(x) = \dfrac{b - a - b\, \ln x}{x^2}.\]

\item En déduire les valeurs des réels $a$ et $b$.
\end{enumerate}

Dans la suite de l'exercice, on admet que la fonction $f$ est définie pour tout réel $x$ de l'intervalle $]0~;~ +\infty[$ par:

\[f(x) = \dfrac{4 + 4\ln x}{x}.\]

\begin{enumerate}[resume]
\item Déterminer les limites de $f$ en $0$ et en $+\infty$.
\item Déterminer le tableau de variations de $f$ sur l'intervalle $]0~;~ +\infty[$. 
\item Démontrer que, pour tout réel $x$ strictement positif, on a :

\[f''(x) = \dfrac{- 4 + 8\ln x}{x^3}.\]

\item Montrer que la courbe $\mathcal{C}_f$ possède un unique point d'inflexion B dont on précisera les coordonnées.
\end{enumerate}

\end{document}
